\chapter[Train Track Case Analysis: Track 8]{Train Track Case Analysis: Track 8}
\label{app:case_analysis_track8}

\section{Case A: Triangles fixed} \label{CaseATrack8}

Let us assume the triangles are fixed by $f_\beta$, hence $\fbd$ and $\fbdb$ both pass over $d$ or $\dbar$ after an even
number of letters. We will call this \emph{Assumption A}. 

\subsection{YOW}\label{YOW}
Suppose $\fbd$ starts at $r$, and $\fbdb$ starts at $p$. In this case, in order to not violate Assumption A, both $\fbd$ and $\fbdb$ must not have any initial twists, and thus
\[
\fbd = r^+ b^+ d...
\]
and 
\[
\fbdb = p^- i^- \dbar...
\]

This is YOW1.

Here the dotted lines must join. Indeed, if $\fbdb$ on the right were to go below $\fbd$ on the left, then it must be the case that $\fbi$ also must go below $\fbd$, otherwise we will be in Situation ???. However, in this case $\fbi$ will necessarily trace:
\[
\fbi = \dbar b^- r^- d p^- i^\times
\]

On the other hand, if $\fbdb$ on the right were to go above $\fbd$ on the left, then similarly $\fbb$ must also go below $\fbdb$. Then in order for $\fbb$ to not trace, we must have
\[
\fbb = d i^+ s^*...
\]
which forces:
\[
\fbd = r^+ b^+ d i^+ s^*...
\]
But now $\fbd$ will violate ???? (odd cross)

Hence the dotted lines must join.

This is YOW2.

here $\fbb$ must go below $\fbi$. Otherwise, $\fbb$ will eventually trace:
\[
\fbb = d p^- i^- \dbar b^\times
\]

Thus we have the situation depicted in YOW3.

Now $\fbo$ is forced to the following:
\[
\fbo = r^+ b^+ d i^+ s^*...
\]

This is YOW3.

Here $\fbo$ can continue with $s^-$ or $s^\circ$. It cannot continue with $s^+$, as that would force it be absorbed:
\[
\fbo = r^+ b^+ d i^+ \hat{s}^+ i^- \dbar b^- r^- \Abso{\fbd}{\fbo}
\]

\subsubsection{Bubo}\label{Bubo}
Suppose $\fbo$ continues with $s^-$. thus we have
\[
\fbo = r^+ b^+ d i^+ \hat{s}^- i^- \dbar b^- r^- d i^*...
\]
This is Bubo1.

Here $\fbp$ cannot start with $i^+$, as it would immediately trace. Thus $\fbp$ either starts with $i^\circ$ or $i^-$.

\subsubsubsection{Gnash}
Suppose
\[
\fbp = i^\circ
\]

Then we have Gnash1.

Here $\fbb$ must end at $p^\circ$. Otherwise, it would trace:
\[
\fbb = d i^+ \hat{p}^+ i^- \dbar b^\times
\]
or
\[
\fbb = d i^+ \hat{p}^+ i^- \dbar r^+ b^\times
\]
or
\[
\fbb = d i^+ \hat{p}^-i^- \dbar b^\times
\]

Thus we must have
\[
\fbb = d i^+ \hat{p}^\circ
\]

This forces
\[
\fbo = ...p^- i^- \dbar b^*...
\]

This is Gnash2.

Here $\fbr$ must end at $b^\circ$. Otherwise, it would trace:
\[
\fbr = b^- r^\times
\]
or
\[
\fbr = b^+ d i^+ p^+ i^- \dbar r^\times
\]

Thus we must have
\[
\fbr = b^\circ
\]

This forces $\fbo$ to continue as shown in Gnash3.

Now consider $\fbs$, we must have
\[
\fbs = i^+ p^- i^- \dbar b^- r^*...
\]

This is Gnash4.

Here $\fbs$ must continue with $r^-$ and to the right of $\fbo$. Otherwise, $\fbo$ will be forced to continue with $r^+$ and to the left of $\fbs$, in which case it would eventually be absorbed:
\[
\fbo = ...r^+ b^+ d i^+ p^+ i^- \dbar r^+ b^+ d i^+ s^+ i^- \dbar b^- r^- \Abso{\fbd}{\fbo}
\]

Thus we must have 
\[
\fbs = i^+ p^- i^- \dbar b^- {r}^-...
\]
but now it will trace, as shown in Gnash5.

This is a contradiction.

\subsubsubsection{Inert}\label{Inert}
Back to the end of Bubo \ref{Bubo}. Suppose
\[
\fbp = i^-...
\]

Consider $\fbi$. It must end here at $o^\circ$, otherwise it would trace:
\[
\fbi = \dbar r^+ b^+ \hat{o}^+ b^- r^- d i^\times
\]
or
\[
\fbi = \dbar r^+ b^+ \hat{o}^- b^- r^- d p^- i^\times
\]

Thus we must have
\[
\fbi = \dbar r^+ b^+ \hat{o}^\circ
\]

This is Inert1.

Here $\fbp$ cannot continue into the path marked with $\times$, as that would lead to it tracing:
\[
\fbp = i^- \dbar r^+ b^+ o^- b^- r^- d p^\times
\]
or 
\[
\fbp = i^- \dbar r^+ b^+ d i^+ s^+ i^- \dbar b^- r^- \Abso{\fbd}{\fbo}
\]

Hence, in the situation depicted in Intert1, $\fbp$ must go into 1 or 2.

\subsubsubsubsection{Callow}
Suppose $\fbp$ goes into 1. 

This is Callow1.

But now notice that $\fbs$ must continue with $i^-$ and to the right of the incoming $\fbo$ to prevent $\fbo$ from being absorbed. But now $\fbs$ will trace:
\[
\fbs = i^- \dbar r^+ b^+ d i^+ s^\times
\]

This is a contradiction.

\subsubsubsubsection{Ennui}\label{Ennui}
Back to the end of Inert \ref{Inert}. Suppose $\fbp$ goes into 2.

This is Ennui1.

Here consider $\fbs$. It must begin with $i^-$ and to the right of $\fbb$, otherwise both $\fbb$ and $\fbo$ will be absorbed between $\fbp$ and $\fbs$.

Thus we must have Ennui2.

Here $\fbb$ is forced to end at $i^\circ$, as either $i^+$ or $i^-$ would cause it to trace immediately after crossing $\dbar$. Thus we have
\[
\fbb = d i^\circ
\]

This is Ennui3. 

Consider $\fbp$. It cannot continue with $b^+$, then it would trace immediately after crossing $d$. So it must either end at $b^\circ$ or continue with $b^-$.

\subsubsubsubsubsection{Brumous}
Suppose it ends at $b^\circ$, so we have
\[
\fbp = i^- \dbar b^\circ
\]

This is Brumous1.

This forces $\fbs$, $\fbo$, and $\fbr$ to continue with $b^+$.

This is Brumous2.

Here $\fbr$ must end at $p^\circ$. Indeed, if it continued with $i^-$ then it would trace:
\[
\fbr = b^+ d i^+ \hat{p}^- i^- \dbar b^- r^\times
\]
While if it continued with $p^+$ it either gets absorbed, if it goes to the left of $\fbs$:
\[
\fbr = b^+ d i^+ \hat{p}^+ i^- \dbar b^- d i^+ \Abso{\fbp}{\fbi}
\]
or trace:
\[
\fbr = b^+ d i^+ \hat{p}^+ i^- \dbar b^- d i^+ \dbar r^\times
\]

Thus we must have 
\[
\fbr = b^+ d i^+ \hat{p}^\circ
\]

This gives us Brumous3.

Here $\fbs$ cannot continue with $r^\circ$ or $r^+$, since in that case $\fbo$ will be forced to continue with $r^+$. But then it will be absorbed:
\[
\fbo = ... r^+ b^+ d i^+ p^+ i^- \dbar b^- d i^+ \Abso{\fbp}{\fbs}
\]

Hence $\fbs$ must continue with $r^-$ and to the right of $\fbo$. But now it will trace:
\[
\fbs = ...r^- b^+ d i^+ p^+ i^- \dbar b^- d i^+ \dbar r^+ b^+ d i^+ s^\times 
\]

This is a contradiction.

\subsubsubsubsubsection{Coterie}
Back to the end of Ennui \ref{Ennui}. Suppose 
\[
\fbp = i^- \dbar \hat{b}^-...
\]

Here $\fbp$ must continue with $b^-$ and to the right of $\fbr$, which forces $\fbr$ to end at $b^\circ$:
\[
\fbr = b^\circ
\]
which in tuen forces $\fbp$ to end as well:
\[
\fbp = i^- \dbar b^- \hat{r}^\circ
\]

This is Coterie1.

Now $\fbs$ and $\fbo$ continue to give Coterie2.

Here $\fbo$ cannot continue with $p^\circ$ or $p^-$, as that would force $\fbs$ to continue with $p^-$ and to the right of $\fbo$, in which case it will trace:
\[
\fbs = ...p^- i^- \dbar b^- r^- b^+ d i^+ \dbar r^+ b^+ d i^+ s^\times
\]

Hence $\fbo$ must continue with $p^+$ and to the left of $\fbs$, but then it will be absorbed:
\[
\fbs = ...p^+ i^- \dbar b^- r^- b^+ d i^+ \Abso{\fbp}{\fbs}
\]

This is a contradiction.

\subsubsection{Gurn}\label{Gurn}
Back to the end of YOW \ref{YOW}. Suppose $\fbo$ ends at $s^\circ$:
\[
\fbo = r^+ b^+ d i^+ s^\circ
\]

This is Gurn1.

Here $\fbp$ can begin and end at $i^\circ$, or it can continue with $i^-$.

\subsubsubsection{Anomy}\label{Anomy}
Suppose
\[
\fbp = i^\circ
\]

This is Anomy1.

Here $\fbb$ must end at $p^\circ$, or it will trace:
\[
\fbb = s i^+ \hat{p}^+ i^- \dbar r^+ b^\times
\]
or
\[
\fbb = s i^+ \hat{p}^- i^- \dbar b^\times
\]

Thus we must have
\[
\fbb = s i^+ \hat{p}^\circ
\]

This is Anomy2.

Here $\fbr$ must begin and end at $b^\circ$:
\[
\fbr = b^\circ
\]

Indeed, if it were to begin with $b^-$ it would immediately trace; if it were to begin with $b^+$ it would either be absorbed, if it goes to the left of $\fbs$:
\[
\fbr = b^+ d i^+ p^+ i^- \Abso{\fbp}{\fbs}
\]
or trace, if it goes to the right of $\fbs$:
\[
\fbr = b^+ d i^+ p^+ i^- \dbar r^\times
\]

This gives us Anomy3.

Consider $\fbi$. It must end here at $o^\circ$. Indeed, otherwise it would trace:
\[
\fbi = \dbar r^+ b^+ \hat{o}^- b^- r^- d p^- i^\times
\]
or
\[
\fbi = \dbar r^+ b^+ \hat{o}^+ b^- r^- d i^\times
\]

Thus we must have 
\[
\fbi = \dbar r^+ b^+ \hat{o}^\circ
\]

Now consider $\fbs$. It must end here at $r^\circ$. Indeed, otherwise it would be absorbed:
\[
\fbs =...r^+ b^+ d i^+ p^+ i^- \Abso{\fbp}{\fbs}
\]
or
\[
\fbs =...r^- b^+ d i^+ p^+ i^- \dbar r^+ b^+ o^- b^- r^- d p^- i^- \dbar b^- r^- \Abso{\fbd}{\fbo}
\]


This gives us Candidate 1, shown in Anomy4.

\subsubsubsection{Bruit}\label{Bruit}
Back to the end of Gurn \ref{Gurn}. Suppose $\fbp$ continues with $i^-$. Again $\fbi$ must end here at $o^\circ$ or it will trace.

This is Bruit1.

Here $\fbp$ must go to the right of $\fbb$. Indeed, if it were to go into the left of $\fbb$, it would either trace:
\[
\fbp = i^- \dbar r^+ b^+ o^- b^- r^- d p^\times
\]
or be absorbed:
\[
\fbp = i^- \dbar r^+ b^+ d i^+ s^+ i^- \dbar b^- r^- \Abso{\fbd}{\fbo}
\]

Thus we must have the situation depicted in Bruit2.

Here $\fbb$ must end at $i^\circ$. Otherwise, it would trace:
\[
\fbb = i^+ \dbar r^+ b^\times
\]
or
\[
\fbb = i^- \dbar b^\times
\]

Thus we must have
\[
\fbb = i^\circ
\]

Here $\fbr$ can begin and end with $b^\circ$, or begin with $b^+$.

\subsubsubsubsection{Diatree}\label{Diatree}
Suppose 
\[
\fbr = b^\circ
\]

This is Diatree1.

Here $\fbp$ must end at $r^\circ$. Indeed, otherwise it would either trace:
\[
\fbp i^- \dbar b^- \hat{r}^+ b^+ d i^+ p^\times
\]
or absorbed:
\[
\fbp i^- \dbar b^- \hat{r}^- b^+ d i^+ \dbar r^+ b^+ d i^+ s^+ i^- \dbar b^- r^- \Abso{\fbd}{\fbi}
\]
or
\[
\fbp i^- \dbar b^- \hat{r}^- b^+ d i^+ \Abso{\fbp}{\fbs}
\]

Thus we must have
\[
\fbp i^- \dbar b^- \hat{r}^\circ
\]

This is Diatree2.

Here $\fbs$ must end at $p^\circ$. Indeed, if it were to continue with $p^+$, it will be absorbed:
\[
\fbs =... p^+ i^- dbar b^- r^- b^+ d i^+ \Abso{\fbp}{\fbs}
\]
While if it were t continue with $p^-$, it is either traces:
\[
\fbs =...p^- i^- \dbar b^- r^- b^+ d i^+ \dbar r^+ b^+ d i^+ s^\times
\]
or is absorbed:
\[
\fbs=...p^- i^- r^- b^+ r^- d i^+ \dbar r^+ b^+ o^- b^- r^- d p^- i^- \dbar b^- r^- \Abso{\fbs}{\fbo}
\]

Thus we must have Diatree3, which is Candidate 2.

\subsubsubsubsection{Fulmine}\label{Fulmine}
Back to the end of Bruit \ref{Bruit}. Suppose $\fbr$ begins with $b^+$. Then it must do so to the left of $\fbp$. Otherwise, it is either absorbed:
\[
\fbr = b^+ d i^+ \Abso{\fbp}{\fbs}
\]
or trace:
\[
\fbr = d i^+ \dbar r ^\times
\]

Thus we must have the situation depicted in Fulmine1.

Here $\fbp$ must end at $b^\circ$. Indeed, if it continued with $b^+$, then it would trace:
\[
\fbp = i^- \dbar \hat{b}^+ d i^+ p^\times
\]
While if it continued with $b^-$, it either is absorbed, if to the left of $\fbs$:
\[
\fbp = i^- \dbar \hat{b}^- d i^+ \Abso{\fbp}{\fbs}
\]
or if to the right of $\fbs$, either it traces:
\[
\fbp = i^- \dbar \hat{b}^- d i^+ \dbar r^+ b^+ o^- b^- r^- d p^\times
\]
or absorbed:
\[
\fbp = i^- \dbar \hat{b}^- d i^+ \dbar r^+ b^+ d i^+ s^+ i^- \dbar b^- r^- \Abso{\fbd}{\fbo}
\]

Thus we must have
\[
\fbp = i^- \dbar \hat{b}^\circ
\]

This gives Fulmine2.

Here $\fbr$ must end at $p^\circ$. Indeed, if it were to continue with $p^-$, it would trace:
\[
\fbr = b^+ d i^+ \hat{p}^- i^- \dbar b^- r^\timed
\]
While if it were to continue with $p^+$, if to the left of $\fbs$, it would be absorbed:
\[
\fbr = b^+ d i^+ \hat{p}^+ i^- \dbar b^- d i^+ \Abso{\fbp}{\fbs}
\]
and if to the right of $\fbs$, it would trace:
\[
\fbr = b^+ d i^+ \hat{p}^+ i^- \dbar b^- d i^+ \dbar r^\times
\]

Thus we must have
\[
\fbr = b^+ d i^+ \hat{p}^\circ
\]

This is Fulmine3.

Here $\fbs$ must end at $r^\circ$. Indeed, if it were to continue with $r^+$, it would be absorbed:
\[
\fbs =...r^+ b^+ d i^+ p^+ i^- \dbar b6- d i^+ \Abso{\fbp}{\fbs}
\]
While if it were to continue with $r^-$, it would either trace:
\[
\fbs =...r^- b^+ d i^+ p^+ i^- \dbar b^- d i^+ \dbar r^+ b^+ d i^+ s^\times 
\]
or be absorbed:
\[
\fbs =...r^- b^+ d i^+ p^+ i^- \dbar b^- d i^+ \dbar r^+ b^+ o^- b^- r^- d p^- i^- \dbar b^- r^- \Abso{\fbd}{\fbo}
\]

Thus we have Candidate 3, shown in Fulmine4.

\subsection{BRR}\label{BRR}
Back to the end of Case A: Triangles fixed \ref{CaseATrack8}. Suppose $\fbd$ starts at $r$, $\fbdb$ starts at $s$. We remind the reader we are still holding Assumption A. 

This is BRR1.

Here the same reasoning as that of YOW \ref{YOW} necessitates that the dotted lines must join. Here $\fbb$ can either go above $\fbi$, or belwo $\fbi$.

Suppose it went above $\fbi$. Then we must have
\[
\fbb = d p^\circ
\]
Otherwise, it would trace if it continued with either $p^-$, $p^+$, or $s^-$. 

This is BRR2.

Here $\fbp$ must begin with $i^-$ and to the right of $\fbs$, otherwise $\fbs$ would be forced to begin with $i^+$ and thus immediately trace.

This is BRR3.

Here $\fbr$ can either begin and end at $b^\circ$, or begin with $b^+$.

\subsubsection{Agog}\label{Agog}
Suppose
\[
\fbr = b^\circ
\]

This is Agog1.

Here $\fbp$ cannot continue with $r^\circ$ or $r^+$, otherwise $\fbi$ will be forced to continue with $r^+$ and to the left of $\fbp$, which would force it to trace:
\[
\fbi = \dbar b^- \hat{r}^+ b^+ d i^\times
\]

Thus we must have
\[
\fbp = i^- \dbar b^- \hat{r}^-...
\]

It cannot continue into the right of $\fbo$, or it will trace after crossing; neither can it go in between $\fbi$ and $\fbo$ for the same reason. But then it can only go into the left of $\fbi$, but then $\fbi$ would still trace.

This is a contradiction.

\subsubsection{Dory}
Back to the end of BRR \ref{BRR}. Suppose $\fbr$ continues with $r^+$.

This is Dory1.

Here $\fbr$ can go into 1,2,3, or 4.

If it were to go into 4 then we are in the same situation as Agog \ref{Agog}, which leads to a contradiction. If it were to go into 3 then it would immediately trace after crossing $d$, as is the case if it were to go into 1. So suppose it goes into 2.

This is Dory2.

Here $\fbs$ is forced to begin with $i^-$ and to the right of the incoming $\fbr$.

Then the situation depicted in Dory3 is forced.

Here $\fbp$ is forced to continue with $r^-$ and to the right of $\fbs$, otherwise $\fbs$ will be forced to continue with $r^+$ and to the left of $\fbp$, in which case it would be absorbed:
\[
\fbs = ...r^+ b^+ d i^+ \Abso{\fbp}{\fbs}
\]

Thus we have Dory4.

Here $\fbp$ must continue with $i^+$ and to the left of $\fbr$. Otherwise, $\fbr$ would be forced to continue with $i^-$ and to the right of $\fbp$, in which case it would trace:
\[
\fbr = b^+ d \hat{i^- \dbar b^- r^\times}
\]

Thus we have Dory5.

Here $\fbo$ must continue with $b^+$ and to the left of $\fbp$. Otherwise, $\fbp$ would be forced to continue with $b^-$ and to the right of $\fbo$, in which case it would trace immediately after crossing.

Thus we have Dory6.

Now $\fbp$ must continue with $b^-$ and to the right of $\fbi$. Otherwise, $\fbi$ would be forced to continue with $b^+$ and to the left of $\fbp$, in which case it would trace immediately after crossing. But now if $\fbp$ continued with $b^-$, it would also trace immediately after crossing.

This is a contradiction.

\subsection{Mien}
Back to the situation depicted in BRR1. We now assume $\fbb$ goes below $\fbi$.

But now $\fbo$ will trace, as shown in Mien1.

\subsection{NUK}
Back to the end of Case A: Triangles fixed \ref{CaseATrack8}. Suppose $\fbd$ starts at $r$, $\fbdb$ starts at $i$. 

The situation is depicted in NUK1.

Here $\fbo$ must go below $\fbp$. Otherwise, we would have
\[
\fbs = \dbar b^- r^- d...
\]
and
\[
\fbp = \dbar b^- r^- d...
\]
in which case at least one of these will trace.

Here $\fbdb$ must come down on $p$ before crossing $\dbar$ in order to adhere to Assumption A (possibly after some number of twisting with $s$).

After crossing, it must go above $\fbb$ otherwise it would be in Situation BAGGED SERPANT. But not it is in Situation ????

This is a contradiction.

\subsection{GAY}\label{GAY}
Back to the end of Case A: Triangles fixed \ref{CaseATrack8}. Suppose $\fbd$ starts at $b$, $\fbdb$ starts at $p$.

Here $\fbr$ can either begin with $o^*$ or $d$.

\subsubsection{Ogee}\label{Ogee}
Suppose 
\[
\fbr = o^*...
\]

This is Ogee1.

Here $\fbd$ must begin with $b^-$ otherwise $\fbo$ would trace.

This is Ogee2.

Here $\fbr$ must end at $o^\circ$ to trace. Here we claim that the dotted lines must join. Indeed, if instead $\fbd$ continues with $r^+$ towards the left of $\fbdb$, then we have the situation shown in Ogee3. But now the dotted lines are in Situation???

On the other hand, in the situatio depicted in Ogee2, if $\fbd$ continues with $r^+$ towards the right of $\fbdb$, then we have Ogee4. Here the dotted lines cannot joint or ??? will be violated. Also, $\fbd$ must contiue with $r^-$ and to the right of $\fbb$ or $\fbb$ will trace.

So now we have Ogee5.

But now $\fbd$ is in Situation ???


Thus in the situation depiated in Ogee2, the dotted lines must join (and $\fbr$ must end). 

This is Ogee6.

Here $\fbo$ can either begin and end at $b^\circ$, or begin with $b^-$.


\subsubsubsection{Chine}\label{Chine}
Suppose
\[
\fbo = b^\circ
\]

This forces
\[
\fbb = r^- d i^+ s^*...
\]

and $\fbi$ must end at $r^\circ$.

This is Chine1.

Here $\fbb$ must end at $s^\circ$. Otherwise it would trace:
\[
\fbb = r^- d i^+ \hat{s}^+ i^- \dbar r^+ b^\times
\]
and
\[
\fbb = r^- d i^+ \hat{s}^- i^- \dbar r^+ o^- b^\times
\]

Thus we must have
\[
\fbb = r^- d i^+ \hat{s}^\circ
\]

This is Chine2.

Here $\fbp$ cannot begin with $i^-$, as it would either trace:
\[
\fbp = i^- \dbar r^+ b^+ o^+ r^- o^- b^- r^- d ^\times
\]
or be absorbed:
\[
\fbp = i^- \dbar r^+ b^+ o^+ r^- d i^+ s^+ i^- \dbar r^+ b^+ \Abso{\fbd}{\fbo}
\]
On the other hand, it cannot begin with $i^+$ as it would immediately trace. Thus we must have
\[
\fbp = i^\circ
\]
This then forces 
\[
\fbs = i^+ p^\circ
\]
This gives Candidate 4. Shown in Chine3.

\subsubsubsection{Droit}\label{Droit}
Back to the end of Ogee \ref{Ogee}. Suppose
\[
\fbo = b^-...
\]

This forces
\[
\fbo = b^- r^- d i^+ s^*...
\]

This is Droit1.

Here $\fbb$ can either begin and end at $r^\circ$ or begin with $r^-$.

\subsubsubsubsubsection{Quaint}\label{Quaint}
Suppose 
\[
\fbb = r^\circ
\]

This is Quaint1.

Here $\fbi$ is forced to end at $b^\circ$, otherwise it would trace:
\[
\fbi = ...b^+ r^- o^- b^- r^- d i^\times
\]
or
\[
\fbi = ...b^- r^- d i^\times
\]
or
\[
\fbi = ...b^- o^- b^- r^- d p^- i^\times
\]

Thus we must have Quaint2.

Now $\fbo$ is forced to end here at $s^\circ$. Otherwise it would either be absorbed:
\[
\fbo = ...s^+ i^- \dbar r^+ b^+ \Abso{\fbd}{\fbo}
\]
or trace:
\[
\fbo = ... s^- i^- \dbar r^+ b^+ r^- o^\times
\]

Thus we must have
\[
\fbo = b^- r^- d i^+ s^\circ
\]

This is Quaint3. 

Now if $\fbp$ were to begin with $i^-$, then it would either trace:
\[
\fbp = i^= \dbar r^+ b^+ o^+ r^+ b^- r^- o^- b^- r^- d p^times
\]
or be absorbed:
\[
\fbp = i^= \dbar r^+ b^+ o^+ r^+ b^- r^- d i^+ s^+ i^- \dbar r^+ b^+ \Abso{\fbd}{\fbo}
\]
On the other hand starting with $i^+$ would immediately trace. Thus we must have
\[
\fbp = i^\circ
\]

This then forces $\fbs$ to begin with $i^+$ and to end there. 

This is Quaint4, which is Candidate5.

\subsubsubsubsection{Abscis}
Back to the end of Droit \ref{Droit}. Suppose
\[
\fbb = r^-...
\]

Then we have Abscis1. here $\fbi$ is forced to end.

But now $\fbo$ cannot continue with $s^\circ$ or $s^+$, otherwise $\fbb$ will be forced to continue with $s^+$ and to the left of $\fbo$ which would lead it to trace:
\[
\fbb = ...s^+ i^- \dbar r^+ b^\times
\]

But then if $\fbo$ were to continue with $s^-$ and to the right of $\fbb$, then it would trace:
\[
\fbo = ... s^- i^- \dbar r^+ o^\times
\]

This is a contradiction.

\subsubsection{Tret}\label{Tret}
Back to the end of GAY \ref{GAY}. Suppose $\fbr$ begins with $d$.

This is Tret1.

Here $\fbd$ must come down at $r$ right before crossing, and with no initial twisting otherwise $o$ will trace.

This is Tret2.

Here $\fbr$ cannot go above $\fbi$, otherwise it would trace:
\[
\fbr = d p^- i^- \dbar b^- r^\times
\]

If $\fbr$ were to go in between $\fbi$ and $\fbdb$, then the dotted lines are in Situation??? shown in Tret3.

So suppose $\fbr$ goes below $\fbdb$. Then we must have
\[
\fbr = d i^+ s^*...
\]
This is Tret4.

In fact it must be
\[
\fbr = d i^+ s^\circ
\]
as otherwise it would trace:
\[
\fbr = d i^+ \hat{s}^+ i^- \dbar r^\times
\]
or
\[
\fbr = d i^+ \hat{s}^- i^- \dbar b^- r^\times
\]

This is Tret5.

Here the dotted lines can either join, or $\fbd$ goes below $\fbdb$.

\subsubsubsection{Fosse}\label{Fosse}
Suppose the dotted lines join. Then this forces
\[
\fbi = \dbar r^+ b^+ o^\circ
\]

This is Fosse1.

Here $\fbb$ must end at $r^\circ$. Otherwise it would trace:
\[
\fbb = r^+ b^\times
\]
or
\[
\fbb = r^- d i^+ p^+ \dbar r^+ b^\times
\]
or
\[
\fbb = r^- d i^+ s^- i^- \dbar b^\times
\]

Thus we must have
\[
\fbb = r^\circ
\]

This is Fosse2.

Here $\fbo$ either begins and ends at $b^\circ$ or begins with $b^-$.

\subsubsubsubsection{Lacuna}\label{Lacuna}
Suppose
\[
\fbo = b^\circ
\]
Then $\fbp$ must begin and end with $i^\circ$. Otherwise it either traces:
\[
\fbp = i^+ p^\times
\]
or is absorbed:
\[
\fbp = i^- \dbar r^+ b^+ d i^+ s^+ i^- \dbar r^+ b^+ \Abso{\fbd}{\fbo}
\]
or
\[
\fbp = i^- \dbar r^+ b^+ o^- b^- r^- d p^- i^- \dbar r^+ b^+ \Abso{\fbd}{\fbo}
\]

Thus we must have
\[
\fbp = i^\circ
\]
This immediately forces 
\[
\fbs = i^+ p^circ
\]

This is Candidate 6, shown in Lacuna1.

\subsubsubsubsection{Exedra}
Back to the end of Fosse \ref{Fosse}. Suppose
\[
\fbo = b^-...
\]
Then the following is forced, shown in Exedra1.

Here $\fbo$ must continue with $i^+$ and to the left of $\fbp$. Otherwise, $\fbs$ will be forced to continue with $i^-$ and to the right of $\fbo$, which would then force it to trace:
\[
\fbs = ...i^- \dbar r^+ b^+ d i^+ s^\times
\]

Then $\fbp$ is forced to begin and end with $i^\circ$. Otherwise, it must begin with $i^-$, then it is either absorbed:
\[
\fbp = i^- \dbar r^+ b^+ d i^+ s^+ i^- \dbar r^+ b^+ \Abso{\fbd}{\fbo}
\]
or trace;
\[
\fbp = i^- \dbar r^+ b^+ b^- b^- r^- d p^\times
\]

Thus we must have
\[
\fbp = i^\circ
\]

This is Exedra2.

Now $\fbo$ cannot continue with $p^\circ$ or $p^-$, as that would force $\fbs$ to continue with $p^-$ and to the right of $\fbo$, in which case it will be trace:
\[
\fbs = i^+ \hat{p}^- i^- \dbar r^+ b^+ d i^+ s^\times
\]

So $\fbo$ must continue with $p^+$ and to the left of $\fbs$, but then it will be absorbed:
\[
\fbo = ...p^+ i^- \Abso{\fbp}{\fbo}
\]

This is a contradiction.

\subsubsubsection{Lemur}
Back to the end of \ref{Tret}. Suppose $\fbd$ goes below $\fbdb$.

This forces the situation depicted in Lemur1. But now the dotted lines are in Situation??

\subsection{OGG}
Back to the end of Case A: Triangles fixed \ref{CaseATrack8}. Suppose $\fbd$ starts at $b$,
$\fbdb$ starts at $s$.

Here again $\fbr$ can either begin with $o^*$ or $d$.

However, the arguments in Ogee \ref{Ogee} and its subsections can be repeated to show that $\fbr$ cannot begin with $o^*$. 

So suppose $\fbr$ begins with $d$. 

This is OGG1.

Here $\fbp$ must begin with $i^-$ and to the right of the incoming $\fbs$. Otherwise, $\fbs$ will be forced to begin with $i^+$ and to the left of $\fbp$, in which case it will immediately trace. Thus we must have 
\[
\fbp = i^- \dbar...
\]

This is OGG2.

Here $\fbr$ cannot go below $\fbdb$, or it will trace:
\[
\fbr = d i^+ s^+ \dbar r^\times
\]

Thus we must have the situation depicted in OGG3.

But now the dotted lines cannot go below $\fbi$ or they will be in Situation ??? BAGGED. On the other hand, going above $\fbi$ puts them in Situation ???.

This is a contradiction.

\subsection{SME}\label{SME}
Back to the end of Case A: Triangles fixed \ref{CaseATrack8}. Suppose $\fbd$ starts at $b$,
$\fbdb$ starts at $i$.

This is SME1.

Here $\fbr$ cannot go above $\fbs$, as at least one of $\fbs$ or $\fbp$ will trace after crossing $\fbd$ again. Here $\fbr$ can continue with $p^*$ or $s^*$.

Suppose it continued with $p^*$. This is SME3. Now it must in fact end there. Otherwise, it would trace:
\[
\fbr = d i^+ \hat{p}^+\dbar r^\times
\]
or
\[
\fbr = d i^+ \hat{p}^+i^- \dbar r^\tmes
\]
or
\[
\fbr = d i^+ \hat{p}^- i^- \dbar b^- r^\times
\]

Thus we must have
\[
\fbr = d i^+ \hat{p}^\circ
\]

This is SME4. But now notice $\fbdb$ will violate Assumption A.

This is contradiction. Thus $\fbr$ must continue with $s^*$. 

This is SME5.

Here $\fbdb$ can either continue with $s^*$ or $p^*$.

\subsubsection{Barm}
Suppose it continued with $s^*$. Then we have the situation depicted in Barm1.

Here $\fbdb$ must continue with $s^+$ , and eventually come down on $p$ before crossing $\dbar$, possibly after some twisting.

This is Barm2.

In Barm3 , the situation is shown with no twisting.

Here the dotted lines can either join, or $\fbdb$ must go below $\fbd$.

Suppose the dotted lines join. This is shown in Barm4. Here $\fbs$ cannot continue with $r^\circ$ or $r^-$, otherwise $\fbp$ would be forced to continue with $r^-$ and to the right of $\fbs$, in which case it traces immediately after crossing $d$. Thus $\fbs$ must continue with $r^+$ and to the left of $\fbb$ (it cannot go in between $\fbb$ and $\fbp$ or it will itself trace eventually after crossing $d$). But now $\fbo$ is forced to begin with $b^-$ and to the right of the incoming $\fbs$.

This is Barm5.

But now $\fbi$ is forced to begin with $p^+$ and to the left of the incoming $\fbo$, and now it will be absorbed:
\[
\fbi = p^+ \dbar r^+ b^+ \Abso{\fbd}{\fbo}
\]

This is a contradiction. Thus the dotted lines cannot join so $\fbdb$ must go above $\fbd$. 

This is Barm6.

But now again $\fbi$ is forced to begin with $p^+$ and to the left of the incoming $\fbdb$, and it will be absorbed.

This is a contradiction.

\subsubsection{Nary}\label{Nary}
Back to the end of SME \ref{SME}. Suppose $\fbdb$ continue with $p^*$.

This is Nary1.

Here $\fbdb$ must continue with $\dbar$ and not to the right of $\fbi$ or it will be in Situation ???. Thus we must have the situation depicted in Nary2.

Here the dotted lines must join. Indeed, if $\fbd$ were to go above $\fbdb$ then $\fbdb$ will be in Situation???. If instead $\fbd$ were to go under $\fbdb$ and above $\fbs$, then $\fbd$ will be in Situation???. If $\fbd$ were to go under $\fbdb$ and also under $\fbp$, then we have the situation depicted in Nary3.

But here the dotted lines will necessarily violate Rule???

Thus the dotted lines must join. This is Nary4.

Here $\fbs$ cannot continue with $r^\circ$ or $r^-$, otherwise $\fbp$ would be forced to continue with $r^-$ and to the right of $\fbs$, in which case it traces immediately after crossing $d$. Thus $\fbs$ must continue with $r^+$.

This is Nary5.

Here $\fbs$ can either go into 1. or 2.

Suppose $\fbs$ were to go into 1. Then we have the situation depicted in Nary6.

Here $\fbp$ must end at $r^\circ$. Otherwise, it is either absorbed:
\[
\fbp = \dbar \hat{r}^+ d \Abso{\fbs}{\fbp}
\]
or trace:
\[
\fbp = \dbar \hat{r}^+ p^\times
\]

This is Nary7.

Here $\fbo$ can either begin and end with $b^\circ$, or begin with $b^-$.

\subsubsubsection{Aulic}\label{Aulic}
Suppose
\[
\fbo = b^\circ
\]
Then we have Aulic1.

Here $\fbs$ can either end at $o^\circ$ or continue with $o^-$.

If it ends here at $o^\circ$, then what is left is $\fbi$ and $\fbb$ twist $k$ times between $s$ and $i$ and end there, giving us Candidate Family 7. Shown in Aulic2.

Suppose instead $\fbs$ continues with $o^-$. Then we have the situation depicted in Aulic3 and Aulic4.

Here if $\fbb$ were to continue with $i^\circ$ or $i^+$, then $\fbs$ continues with $i^+$ and to the left of $\fbb$. But then it will be absorbed:
\[
\fbs =...i6+ p^+ \dbar r^+ b^+ \Abso{\fbd}{\fbo}
\]

On the other hand, if $\fbb$ were to continue with $i^-$ and to the right of $\fbs$, then it will trace:
\[
\fbb = ...i^- p^+ \dbar r^+ d i^+ p^+ \dbar r^+ b^\times
\]

This is a contradiction.

\subsubsubsection{Oread}\label{Oread}
Back to the end of Nary \ref{Nary}. Suppose
\[
\fbo = b^-...
\]
This is Oread1.

Here $\fbo$ must continue with $p^-$ and to the right of $\fbi$. Otherwise, $\fbb$ will trace.

This is Oread2.

Here $\fbi$ can either begin and end here at $p^\circ$, or begin with $p^+$. Suppose 
\[
\fbi = p^\circ
\]
This is Oread3. 

Here $\fbo$ must continue with $i^-$ and to the right of $\fbb$. Otherwise $\fbb$ must continue with $i^+$, and it will trace.

But now either $\fbo$ traces:
\[
\fbo =...i^-p^+ \dbar r^+ d i^+ p^+ \dbar r^+ b^+ o^\times
\]
or it is absorbed:
\[
\fbo = ...i^- p^+ \dbar r^+ d i^+ s^- i^- \dbar b^- r^- d p^- i^- \dbar r^- d \Abso{\fbs}{\fbp}
\]

Thus we must have
\[
\fbi = p^+...
\]

This is Oread4.

This eventually forces the situation depicted in Oread5, and then Oread6 which is Candidate 8.

\section{Case B,C,D: Triangles Swapped}\label{CaseBCDTrack8}
Now we assume the triangles are swapped.

We first assume that $\fbd$ passes $\dbar$ after an odd number of letters. 

\subsection{FUN}\label{FUN}
Supposing $\fbd$ passes $\dbar$ after an odd number of letters. Suppose $\fbd$ starts at $p$, and suppose $\fbdb$ starts at $d$. Then after possibly twisting with $i$, $\fbd$ must come down on $p^+$ immediately before crossing. 

This is FUN1.

Here $\fbp$ and $\fbs$ must begin with $b^+$ and to the left of $\fbr$. Otherwise, $\fbr$ will be forced to continue with $b^-$ and immediately trace. 

This is FUN2.

Now we claim that in fact $\fbd$ cannot have initial twists. Indeed, if it did then $\fbb$ would be forced to do the following:
\[
\fbb = i^+... p^+ \dbar r^*...
\]
and here it must continue with $r^\circ$ or $r^-$ to not immediately trace. But then this forces
\[
\fbi = r^- d p^- i^\times
\]

Thus there cannot be initial twists.

This is FUN3.

Consider the dotted lines $\fbd$ and $\fbdb$. Here they can either join, or $\fbd$ must go below $\fbdb$ in order to not put $\fbdb$ in Situation ???. If it were the latter case, then we would have the situation depicted in FUN4.

But now, the dotted lines are in Situation??? Thus the dotted lines must join. This is FUN5. 

Here $\fbo$ must continue with $r^+$ and to the left of $\fbi$, or $\fbi$ will be forced to begin with $r^-$ and to the right of $\fbo$, in which case it will trace:
\[
\fbi = r^- d p^- i^\times
\]

Thus the situation depicted in FUN6 is forced.

Now $\fbb$ must begin with $i^-$ and to the right of both $\fbp$ and $\fbs$. Otherwise $\fbp$ and $\fbs$ will be forced to continue with $i^+$ and thus at least one of them will then immediately trace. But now $\fbb$ will trace:
\[
\fbb = i^- \dbar b^\times
\]

This is a contradiction.

\subsection{XOY}
Supposing $\fbd$ passes $\dbar$ after an odd number of letters. Suppose $\fbd$ starts at $p$, and suppose $\fbdb$ starts at $o$. 

This case is impossible as $\fbr$ would immediately trace:
\[
\fbr = \dbar r^\times
\]

\subsection{FAT}\label{FAT}
Supposing $\fbd$ passes $\dbar$ after an odd number of letters. Suppose $\fbd$ starts at $p$, and suppose $\fbdb$ starts at $b$.

This is Fat1.

Here $\fbd$ must begin with $b^+$ or $\fbr$ will trace. Similarly $\fbdb$ must continue with $r^+$ or $\fbp$ will trace.

Thus we have FAT2.

Here $\fbs$ must begin with $r^+$ and to the left of $\fbp$ or $\fbp$ will trace.

This gives FAT3.

Then the situation in FAT4 is forced. Here because $\fbp$ cannot escape the left hand side, it twists between $r$ and $o$ $k$ times and end there. Then $\fbs$ is forced to end at whichever $r$ or $o$ is left. Then $\fbi$ is forced to end at $p^\circ$.

This is Candidate Family 9.

\subsection{KNH}
Supposing $\fbd$ passes $\dbar$ after an odd number of letters. Suppose $\fbd$ starts at $p$, and suppose $\fbdb$ starts at $r$.

Then we must have the situation in KNH1, where $\fbp$ must go below $\fbr$ to not trace. But now $\fbr$ must trace.

This is a contradiction.

\subsection{MZG}
Supposing $\fbd$ passes $\dbar$ after an odd number of letters. Suppose $\fbd$ starts at $s$, and suppose $\fbdb$ starts at $d$.

This is MZG1. 

Here $\fbp$ and $\fbs$ must begin with $b^+$ and to the left of $\fbr$ or $\fbr$ will trace.

Then we have MZG2.

Here $\fbp$ can continue with $d$ or $o^*$.

Suppose $\fbp$ were to continue with $o^*$. Then this is MZG3. Here the dotted lines must join. Indeed, if not the only possibility is $\fbd$ goes below $\fbdb$, otherwise $\fbdb$ will be in Situation???. But then we would have
\[
\fbd = ... d b^+ o^- b^- r^- d...
\]
which would then put the dotted lines in Situation???

So we must have the situation depicted in MZG3.

Here let us assume that $\fbd$ has no initial twists. Then we have the situation depicted in MZG4.

Here $\fbp$ must continue with $o^+$ and to the left of $\fbr$, otherwise $\fbr$ will be absorbed between $\fbp$ and $\fbs$. This gives us MZG5.

Here we claim that $\fbp$ cannot continue with $b^-$. Indeed, if it did then we would be in the situation depicted in MZG6. But here $\fbp$, $\fbi$, and $\fbo$ are stuck in a path mill ????.

Thus in the situation depicted in MZG5, we must have $\fbp$ continue with $b^\circ$ or $b^+$, thus $\fbs$ must continue with $b^+$. But now $\fbr$, $\fbp$, and $\fbs$ are stuck in a counterclockwise path mill???


Thus we must assume that $\fbd$ has initial twists. This forces $\fbb$ to cross $d$ after twisting. But this then forces the situation depicted in MZG7.

Here notice $\fbd$ is in Situation???

This is a contradiction.


\subsection{GUM}
Supposing $\fbd$ passes $\dbar$ after an odd number of letters. Suppose $\fbd$ starts at $s$, and suppose $\fbdb$ starts at $o$.

This case cannot occur as $\fbr$ will immediately trace after crossing.

\subsection{WAK}
Supposing $\fbd$ passes $\dbar$ after an odd number of letters. Suppose $\fbd$ starts at $s$, and suppose $\fbdb$ starts at $b$.

This is WAK1.

Here $\fbdb$ must begin with $b^+$ otherwise $\fbr$ will trace. Also here $\fbi$ can either go above $\fbd$ or below $\fbd$. 

Suppose it were to go above. Then WAK2 is forced. But now $\fbd$ must have initial twists, otherwise $\fbo$ will eventually trace. This gives WAK3. But now $\fbb$ will eventually trace after crossing:
\[
\fbb = i^+...p^+ \dbar b^\times
\]

Thus $\fbi$ must go below $\fbd$. This gives WAK4. Here if the dotted lines continue with $d$, then they would be in Situation ???, so they must continue with $o^*$. But now notice by the same reasoning as before $\fbd$ must have initial twists otherwise $\fbo$ will trace after crossing. But then again $\fbo$ will eventually trace, again by the same reasoning as before.

This is a contradiction.

\subsection{MAU}
Supposing $\fbd$ passes $\dbar$ after an odd number of letters. Suppose $\fbd$ starts at $s$, and suppose $\fbdb$ starts at $r$.

This is MAU1.

But here at least one of $\fbp$ and $\fbs$ will immediately trace after crossing.

This is a contradiction.

\subsection{EEN}
Supposing $\fbd$ passes $\dbar$ after an odd number of letters. Suppose $\fbd$ starts at $i$, and suppose $\fbdb$ starts at $d$. This is EEN1, where $\fbd$ can possible have initial twists. Here $\fbi$ must begin with $r^-$ and to the right of $\fbb$, or $\fbb$ will trace. This gives EEN2.

Here $\fbi$ must continue with $p$ or $s$. Suppose $\fbi$ contiues with $p$. This is EEN3.

Now if $\fbd$ were to have initial twists, then $\fbr$ will be forced to follow the path as shown in EEN4. This then forces $\fbp$ to begin with 
\[
\fbp = b^+ o^*...
\]
as shown.

But now $\fbp$, $\fbs$, and $\fbr$ are in a three-way path mill.

So $\fbd$ cannot have initial twists. This is EEN5. Here $\fbo$ can either end at $i^\circ$, or begin with $i^-$. Suppose it were to begin with $i^-$. This is EEN6. Then it continues with 
\[
\fbo = i^- \dbar r^- d...
\]
and here if the next letter were $p$, then $\fbo$, $\fbi$, and $\fbb$ will be in a three-way clockwise path mill. Thus the next letter must be $s^-$. So we have
\[
\fbo = i^- \dbar r^- d \hat{s}^- i^- \dbar...
\]
This forces
\[
\fbp = b^+ o^*...
\]

This is EEN7. Now here $\fbp$, $\fbs$, and $\fbo$ are stuck in a three-way couterclockwise path mill.

This is a contradiction. 

Thus going back to the situation depicted in EEN2, $\fbi$ must continue with $s^*$. This is EEN8.

But now $\fbi$ must continue with $p^-$ and to the right of $\fbr$, or $\fbr$ will trace immediately after crossing $\dbar$. But now $\fbi$ traces. This is a contradiction.


This concludes the situation where $\fbd$ passes an odd number of letters before crossing. Now we assume $\fbd$, Now, suppose instead that $\fbd$ passes over $\dbar$ after an even number of letters. From previous arguments, we can then
conclude that $\fbdb$ passes over $\dbar$ after an odd number of letters. Then the only possibility for $\fbdb$ is to start at $b$, $r$, or $o$. While for $\fbd$ the only possibility is to start at $p$, $s$, or $i$.

\subsection{MYU}\label{MYU}
Supposing both $\fbd$ passes $\fbar$ after an even number of letters. Suppose $\fbd$ starts at $p$, and suppose $\fbdb$ starts at $b$.

This is MYU1. Here $\fbr$ must go below $\fbi$ or it will immediately trace. But now $\fbi$ will trace. As shown in MYU2.

This is a contradiction.

\subsection{QBC}\label{QBC}
Supposing both $\fbd$ passes $\fbar$ after an even number of letters. Suppose $\fbd$ starts at $p$, and suppose $\fbdb$ starts at $o$. This is QBC1. But now $\fbr$ will immediately trace after crossing.

This is a contradiction.

\subsection{OWD}\label{OWD}
Supposing both $\fbd$ passes $\fbar$ after an even number of letters. Suppose $\fbd$ starts at $p$, and suppose $\fbdb$ starts at $r$.

Here again $\fbr$ must continue with $b$ after crossing. But this forces $\fbp$ to trace, as shown in OWD1.

This is a contradiction.

\subsection{BFI}\label{BFI}
Supposing both $\fbd$ passes $\fbar$ after an even number of letters. Suppose $\fbd$ starts at $p$, and suppose $\fbdb$ starts at $r$.

This is a contadiction.

Now suppose $\fbd$ begins at $s$. The arguments of MYU \ref{MYU}, QBC \ref{QBC}, OWD \ref{OWD}, and BFI \ref{BFI} can all be repeated to rule this scenario out.

\subsection{CQD}
Supposing both $\fbd$ passes $\fbar$ after an even number of letters. Suppose $\fbd$ starts at $i$, and suppose $\fbdb$ starts at $o$.

Then this is CQD1.

Here $\fbi$ is forced to begin with $r^-$ and to the right of $\fbb$, or $\fbb$ will trace. This is CQD2. Now notice that the next letter of $\fbd$ must either be $s$ or $p$, but in the latter case $\fbi$ will trace. So we must have
\[
\fbd = i^+ s^- \dbar...
\]
and
\[
\fbi = r^- d p^*...
\]

But now $\fbd$ continues with
\[
\fbd = i^+ s^- \dbar r^+ o^...
\]
This forces $\fbo$ to trace. As shown in CQD3.

This is a contradiction.

\subsection{KYD}
Supposing both $\fbd$ passes $\fbar$ after an even number of letters. Suppose $\fbd$ starts at $i$, and suppose $\fbdb$ starts at $b$.

This is KYD1.

Here $\fbp$ and $\fbs$ must begin with $r^-$ and to the right of $\fbb$. Otherwise $\fbb$ will immediately trace. But then notice at least one of $\fbp$ or $\fbs$ will immediately trace after crossing $d$.

This is a contradiction.

\subsection{FHR}\label{FHR}
Supposing both $\fbd$ passes $\fbar$ after an even number of letters. Suppose $\fbd$ starts at $i$, and suppose $\fbdb$ starts at $r$. 

This is FHR1.

Here we either have
\[
\fbd = i^+ s^+...
\]
or 
\[
\fbd = i^+ p^+...
\]

\subsubsection{Eyas}
Suppose
\[
\fbd = i^+ s^+...
\]
Then the situation in Eyas1 is forced. The dotted lines are forced to connect since otherwise either $\fbdb$ enters Situation $\beta$, or the dotted lines enter Situation $\alpha$.

Here if $\fbp$ continued with $o^+$ and to the left of $\fbs$, then $\fbp$, $\fbb$, and $\fbs$ will be in a three-way couterclockwise path mill as shown in Eyas2.

So we must have $\fbs$ continue with $o^-$, and $\fbp$ ending at $o^\circ$. 

This is Eyas3.

Here we claim that $\fbi$ cannot continue with $p^+$ and to the left of $\fbs$. Indeed, if so then we would have the situation depicted in Eyas4. But here $\fbi$ must continue with $s^-$ or $\fbr$ will trace, but then $\fbi$ immediately traces as well.

This gives Eyas5. But here notice that $\fbo$, $\fbo$, and $\fbi$ are in a three-way clockwise path mill. 

This is a contradiction,

\subsubsection{Fard}\label{Fard}
Back to the end of FHR \ref{FHR}. Suppose
\[
\fbd = i^+ p^+...
\]

This is Fard1.

Here the dotted lines must connect. Otherwise, if $\fbdb$ went under $\fbd$, then $\fbdb$ will be in Situation $\beta$. On the other hand, if $\fbd$ went below $\fbdb$, then we would have Fard2, in which case the dotted lines are in Situation $\alpha$.


Thus we must have Fard3.

Here $\fbp$ can either begin with $d$, or begin with $o$. Suppose
\[
\fbp = d...
\]

Then $\fbs$ must begin with $o$. This is Fard4.

Here $\fbo$ cannot go up to $s$ because then if $\fbp$ continued with $s^\circ$ or $s^-$, $\fbo$ will be forced to continue with $s^-$, but then it will be absorbed between $\fbs$ and $\fbp$. On the other hand, if $\fbp$ continued with $s^+$, then it will be absorbed between $\fbd$ and $\fbo$.

Thus we must have the situation depicted in Fard5. But now $\fbi$ must begin with $b^-$ and to the right of $\fbo$, otherwise $\fbo$ will immediately trace. But now $\fbi$ will trace:
\[
\fbi = b^- r^- d p^- i^\times
\]

This is a contradiction.

So suppose
\[
\fbp = o^*...
\]

This gives the situation depicted in Fard6.

Here $\fbi$ can either begin and end with $b^\circ$, or begin with $b^+$.

\subsubsubsection{Dwaal}
Suppose
\[
\fbi = b^\circ
\]

This gives Dwaal1.

Here $\fbs$ must continue with $r^-$ and to the right of $\fbb$. Otherwise, $\fbb$, $\fbs$, and $\fbp$ will be in a three-way counterclockwise path mill.

Thus we have Dwaal2.

Here $\fbr$ and $\fbo$ twist among themselves between $p$ and $i$ and end there, since $\fbr$ can never escape or it will trace. This then forces $\fbo$ to end here at $s^\circ$, and $\fbb$ to end here at $r^\circ$. 

This gives Cnadidate Family 10.

\subsubsubsection{Fremd}
Back to the end of Fard \ref{Fard}. Suppose
\[
\fbi = b^+...
\]
This is Fremd1.

Here $\fbi$ must continue with $o^+$ and to the left of $\fbp$ and $\fbs$, otherwise $\fbp$ will trace after crossing $d$. This is Fremd2.

Here $\fbb$ must continue with $r^+$ and to the left of $\fbi$, otherwise $\fbi$ will continue with $d$ and then $p$. Then it must continue with $p^-$, or $\fbr$ will trace, but then $\fbi$ itself traces. So we have the situation depicted in Fremd3. Here $\fbp$ and $\fbs$ must continue with $o^+$ and to the left of $\fbb$, otherwise $\fbb$ will immediately trace. So we have Fremd4. This continues and we see that $\fbp$, $\fbs$, $\fbb$, and $\fbi$ are in a four-way counterclockwise path mill.

This is a contradiction.

